% Generated draft: requires mathematical and source review.
\documentclass[11pt]{article}
\usepackage[margin=25mm]{geometry}
\usepackage{fontspec}
\setmainfont{Latin Modern Roman}
\usepackage{amsmath,amssymb,mathtools,unicode-math}
\setmathfont{Latin Modern Math}
\usepackage{enumitem,needspace,xurl,hyperref}
\setlength{\parindent}{0pt}
\setlength{\parskip}{5pt}
\setlength{\emergencystretch}{3em}
\newcommand{\N}{\mathbb{N}}
\newcommand{\Z}{\mathbb{Z}}
\newcommand{\Q}{\mathbb{Q}}
\newcommand{\R}{\mathbb{R}}
\newcommand{\C}{\mathbb{C}}
\newcommand{\F}{\mathbb{F}}
\newcommand{\A}{\mathbb{A}}
\newcommand{\PP}{\mathbb{P}}
\newcommand{\E}{\mathbb{E}}
\newcommand{\eps}{\varepsilon}
\newcommand{\dd}{\,\mathrm{d}}
\begin{document}
\section*{UnsolvedMath ID: 30005218}
\textbf{Source problem code:} catalogue entry
\subsection{Definitions and mathematical statement}
The supplied question concerns affine type \(A\) and uses the symbols \(N_{\lambda_{i+1}}\) and \(N_{\mu_i}\), without defining these objects, their ambient algebra or module, their coefficient ring, or their indexing sets. The subscripts below reproduce the supplied cleaned statement; their intended mathematical meaning is unresolved.

In general, given a coefficient ring \(R\), a designated set \(R_{\geq 0}\subseteq R\) of nonnegative coefficients, and a specified expansion
\[
N_{\lambda_{i+1}}=\sum_{\mu}c_{\lambda,\mu,i}N_{\mu_i},
\]
positivity of that expansion means
\[
c_{\lambda,\mu,i}\in R_{\geq 0}
\quad\text{for every index \(\mu\) occurring in the expansion.}
\]
This is a conditional explanation of the word ``positive,'' not a specification of the missing coefficient ring, index ranges, basis properties, or permitted sums.

\textbf{Source question.} In affine type \(A\), is the decomposition of \(N_{\lambda_{i+1}}\) in terms of the \(N_{\mu_i}\) positive?

A fully quantified conjecture cannot be recovered from this record: the affine rank, admissible values of \(i,\lambda,\mu\), meaning of the indexed \(N\)-objects, and intended notion of coefficient positivity are not supplied. In particular, the record does not determine whether \(i\) denotes a level of a family of pre-canonical bases or forms part of a weight subscript. The supplied literature assessment associates the question with pre-canonical bases, but does not provide their definitions. The corrupted original excerpt also mentions a second problem concerning positivity properties of Kazhdan--Lusztig polynomials and then reports an Elias--Williamson positivity result. Those passages do not specify an additional unresolved mathematical assertion and cannot be used to identify the displayed conjecture with ordinary Kazhdan--Lusztig polynomial positivity.
\subsection{Short English statement}
The record asks whether certain expansion coefficients in affine type \(A\) are nonnegative. It does not define the elements, indices, coefficient ring, or meaning of nonnegativity, so the question remains a statement template requiring source clarification and human mathematical review. The supplied dated assessment distinguishes this pre-canonical-basis question from the reported positivity theorem for Kazhdan--Lusztig polynomials.
\subsection{Sources}
\begin{itemize}
\item Supplied local registry record: permanent UnsolvedMath ID \(30005218\), title ``Positivity in Affine Kazhdan--Lusztig Theory,'' exact source problem code \texttt{catalogue entry}. The registry supplies the labels \texttt{partially\_solved} and \texttt{corrected\_verified}; these are source labels, not independent verification of the mathematics.
\item Supplied report citation: \emph{Character Theory and Categorification (2023)}, Oberwolfach Reports, pp.\ 2267--2333, DOI \url{https://doi.org/10.4171/owr/2022/39}. The registry separately gives proposed year \(2022\). The supplied description says that Oberwolfach Reports document mathematical workshops and their open-problem discussions.
\item Attribution appearing in the original excerpt: ``Some aspects of Kazhdan--Lusztig polynomials,'' Nicolas Libedinsky, joint work with Leonardo Patimo, David Plaza and Geordie Williamson. The record does not separately identify a proposer of this conjecture. It attributes the reported Kazhdan--Lusztig positivity breakthrough to Ben Elias and Geordie Williamson.
\item Supplied literature reference: Patimo--Plaza--Williamson, pre-canonical bases, \url{https://doi.org/10.1016/j.aim.2022.108255}.
\item Supplied literature reference: Elias--Williamson, Soergel Hodge theory, \url{https://doi.org/10.4007/annals.2014.180.3.6}.
\item Supplied literature reference: original Oberwolfach report, \url{https://doi.org/10.4171/owr/2022/39}.
\item Literature assessment dated exactly \texttt{2026-08-22}: the supplied curation labels the problem partially solved and states that affine type \(A\) pre-canonical-basis positivity remains open beyond verified ranges, while Elias and Williamson proved positivity of Kazhdan--Lusztig polynomials in full generality. The verified ranges are not identified. This is supplied dated curation derived from the supplied OWR material, not an independently verified literature assessment; its time-sensitive status and cited claims require independent checking.
\end{itemize}
\end{document}
